\lecture{10}{24. September 2025}{Internal Flows: Fully Developed Laminar Flows}

\section{Internal Incompressible Viscous Flow}

\subsection{Internal Flow Characteristics}

\subsubsection{Laminar Versus Turbulent Flow}
The pipe flow regime, i.e. if the flow is turbulent or laminar, is determined by the Reynolds number, briefly mentioned in \autoref{sec:sigdim}, $\mathrm{Re} = \frac{\rho \overline{V} D}{\mu}$. Under normal conditions, the transition to turbulence occurs at roughly $\mathrm{Re} \approx \num{2300} $ for flow in pipes. For water in a $Ø=\qty{1}{in.}$ pipe, this corresponds to an average speed of $V \approx \qty{0,1}{\frac{m}{s}}$. Turbulence occurs when the viscous forces are not able to damp out the random fluctuations in fluid motion. This means that a fluid of higher viscosity (e.g. motor oil compared to water) is able to be flowing at a higher rate before becoming turbulent. On the other hand, high-density fluids will exacerbate the inertial forces due to random fluctuations in motion, and therefore this fluid will become turbulent at a lower flow rate.

\subsubsection{The Entrance Region}

